\begin{figure}[ht]
\centering
\begin{tikzpicture}[x=1cm, y=1cm, font=\small]

% Two-panel: p-value histogram under the null (flat) and under a
% well-powered alternative (right-skewed, mass at small p). Bins on
% [0, 1] in 0.05-wide buckets, ten bins shown.

% ---- Panel A: under H0 (uniform) ----
\begin{scope}[shift={(0,0)}]
  \draw[->, thick] (-0.1, 0) -- (5.5, 0)
    node[below, font=\scriptsize] {$p$-value};
  \draw[->, thick] (0, -0.1) -- (0, 3.4)
    node[left, font=\scriptsize] {frequency};
  \foreach \k in {0, 1, 2, 3, 4, 5, 6, 7, 8, 9} {
    \pgfmathsetmacro\xpos{\k*0.5}
    \pgfmathsetmacro\height{1.4 + 0.20*rand}
    \draw[fill=engblue!50, draw=engblue]
      (\xpos, 0) rectangle (\xpos+0.5, \height);
  }
  \draw[dashed, engred, thick] (0, 1.4) -- (5.0, 1.4);
  \node[font=\scriptsize, engred, anchor=west] at (3.2, 1.6)
    {expected (uniform)};
  \node[anchor=south, font=\small\bfseries] at (2.5, 3.5)
    {(a) Under the null};
  \foreach \k/\lab in {0/0, 5/0.5, 10/1.0} {
    \pgfmathsetmacro\xpos{\k*0.5}
    \node[anchor=north, font=\scriptsize] at (\xpos, -0.05) {\lab};
  }
\end{scope}

% ---- Panel B: under a well-powered alternative (mass at small p) ----
\begin{scope}[shift={(7.0,0)}]
  \draw[->, thick] (-0.1, 0) -- (5.5, 0)
    node[below, font=\scriptsize] {$p$-value};
  \draw[->, thick] (0, -0.1) -- (0, 3.4)
    node[left, font=\scriptsize] {frequency};
  \foreach \k/\h in {0/3.1, 1/1.6, 2/1.0, 3/0.7, 4/0.55, 5/0.45,
                     6/0.40, 7/0.35, 8/0.30, 9/0.25} {
    \pgfmathsetmacro\xpos{\k*0.5}
    \draw[fill=engred!50, draw=engred]
      (\xpos, 0) rectangle (\xpos+0.5, \h);
  }
  \node[anchor=south, font=\small\bfseries] at (2.5, 3.5)
    {(b) Under the alternative};
  \foreach \k/\lab in {0/0, 5/0.5, 10/1.0} {
    \pgfmathsetmacro\xpos{\k*0.5}
    \node[anchor=north, font=\scriptsize] at (\xpos, -0.05) {\lab};
  }
\end{scope}

\end{tikzpicture}
\caption[Distribution of $p$-values under null and alternative]{The
$p$-value's own distribution is the fastest way to see what a
$p$-value is. Under the null, $p$ is uniform on $[0, 1]$ (panel a):
every twenty repetitions of a true-null experiment produces, on
average, one $p < 0.05$ by chance. Under a well-powered alternative
(panel b), the distribution piles up near zero; the height of the
leftmost bin is the power of the test. A literature whose $p$-value
distribution shows a sharp spike just below $0.05$ but flat above
is exhibiting the signature of publication bias and selective
reporting.}
\label{fig:vol01:ch08:p-value-distributions}
\end{figure}
